\subsection{An obstruction defined before its vanishing}

\begin{statement}{A.1}{Proved}{The relative cohomology object of the two signs}
Assume the two normalized frame characters of supplied R4,
$a_3=\epsilon_3T$ and $a_2=\epsilon_2W:\HH\to\Ftwo$.
Let $I=[0,1]$, with its endpoints labelled $3,2$, and put
\[
 O=H^1(I,\partial I;\Ftwo)
   =\Ftwo^{\{3,2\}}/\Ftwo(1,1).
\]
The continuous map
$o_{\HH}(h)=[(a_3(h),a_2(h))]$ is a homomorphism.
It is the obstruction to extending the specified endpoint signs to a
locally constant sign on $I$. This definition does not assume that it
vanishes, and the R2 element has nonzero obstruction.
\end{statement}
\begin{proof}
The long exact sequence of the pair, or the cellular complex of an
interval, gives
\[
0\longrightarrow\Ftwo\xrightarrow{t\mapsto(t,t)}
\Ftwo^2\xrightarrow{(u,v)\mapsto u+v}\Ftwo\longrightarrow0.
\]
A locally constant function on the connected interval is constant;
its boundary values therefore agree. Conversely equal values extend
uniquely as a function, after the values have been specified. The
connecting homomorphism is their class modulo the diagonal. Composing
with the continuous homomorphism $(a_3,a_2)$ proves the assertions.
The interval is a precise algebraic model for the two-sign gluing
problem, not an asserted geometric path in a moduli stack.
\end{proof}

\begin{statement}{A.17}{Refuted}{The standard genus-two matrix action is not an action of the five-point sphere group}
The usual symplectic representation of $B_5$ from the genus-two
double cover does not descend even projectively to
$\Gamma_{0,[5]}$, including at level twelve. It therefore
cannot be substituted without a new comparison for the
five-point frame equations.
\end{statement}
\begin{proof}
Take a chain basis $e_1,\ldots,e_4$ of homology with
$J_{i,i+1}=1$, $J_{i+1,i}=-1$, and other entries zero.
It is unimodular. Its transvections
$T_i=I+e_i e_i^{\mathsf T}J$ satisfy the Artin braid
and distant-commutation relations. Direct multiplication gives
\[
 (T_1T_2T_3T_4)^5=-I,
\]
but the sphere relator is
\[
 T_4T_3T_2T_1^2T_2T_3T_4
 =-I+2(e_1+e_3)e_4^{\mathsf T}.
\]
It is not scalar modulo three, four or twelve. Killing
the central full twist thus does not kill the sphere
relator, and passing to the projective symplectic group
does not fix the failure. All matrix identities are
checked in \texttt{check\_genus\_two\_descent.py}.
Geometrically the curve $y^2=x^5-1$ has six branch
points, the five roots and infinity. A double cover
cannot have exactly five simple branch points, by
Riemann--Hurwitz. The extra point must be retained in
any proposed functorial comparison with five-point data.
This refutes the immediate substitution, not the
possibility of a different comparison theorem.
\end{proof}

\begin{statement}{A.2}{Proved}{Descent to every multiplier and the cancelling sign}
Write $G=\GZ$, $\Gamma=G_{\Q}$, and $K=\Gamma\cap\HH$.
For $g\in G$, choose $\sigma\in\Gamma$ with
$\chi(\sigma)=\lambda(g)$ and define
\[
 o(g)=[(a_3(g\sigma^{-1}),a_2(g\sigma^{-1}))]\in O.
\]
This is well defined, continuous and a homomorphism. Identifying $O$
with $\{\pm1\}$ by $[(u,v)]\mapsto(-1)^{u+v}$, its inverse is the
unique sign that cancels it. Its additive exponent is exactly R4's
$\Delta(g)$.
\end{statement}
\begin{proof}
The cyclotomic character is onto, so a choice exists. Replacing
$\sigma$ by $\sigma'$ changes $h=g\sigma^{-1}$ to
$h(\sigma\sigma'^{-1})$, where the second factor belongs to $K$.
R4 and the common-field sign comparison say that the two characters
agree on $K$; the change is therefore diagonal and dies in $O$.
For $g_i=h_i\sigma_i$ the product is
$h_1(\sigma_1h_2\sigma_1^{-1})\sigma_1\sigma_2$.
The conjugation invariance of $a_3+a_2$ established in R4 gives the
homomorphism law, in this order of composition; reversing the order
gives the same additive value. The map $\HH\times\Gamma\to G$ is a
continuous compact surjection. The continuous function just defined
is constant on its fibres and hence descends continuously.
Finally $\{\pm1\}$ has exponent two, so the inverse sign is the same
sign. The definition agrees with
$\Delta(g)=a_3(g\sigma^{-1})+a_2(g\sigma^{-1})$, without a claim
that either sign is $+1$. The cancellation identity
$(-1)^{\Delta(g)}o(g)=1$ is a definition; the assertion $o(g)=1$
on $\GT$ is an additional theorem still required.
\end{proof}

\begin{statement}{A.3}{Refuted}{A vanishing class does not select a gauge-invariant torsor trivialization}
The gauge-invariant two-sign datum does not canonically select one
of its two lifts, even when $\Delta=0$. On the normalized subgroup
$\HH$, the ordered pair of signs is a finer continuous invariant.
It has no extension to an abelian-valued continuous homomorphism on
$G$ retaining both components.
\end{statement}
\begin{proof}
For $\omega\in O$ its inverse image
$P_\omega\subset\Ftwo^2$ is a torsor under the diagonal $\Ftwo$.
The simultaneous change $(u,v)\mapsto(u+1,v+1)$ preserves the
class and interchanges its two lifts. A choice natural under that
gauge transformation would be a fixed point, which does not exist.
In particular $P_0=\{(0,0),(1,1)\}$ has two trivializations as a
torsor. Declaring $(0,0)$ preferred uses a chosen signed coordinate;
it is not invariant under the allowed overall-sign change.
An ordinary torsor over a point is always trivializable; calling
$P_1$ a nontrivial torsor over a point would be incorrect. The
obstruction is the class in the quotient, or equivalently the
equivariant sign representation with character $\Delta$.

There exists $k\in K$ with $(a_3(k),a_2(k))=(1,1)$: the extension
$\Q^{\rm ab}((3/4)^{1/4})/\Q^{\rm ab}$ is quadratic, as shown in
A.12 below, and a nontrivial automorphism extends to an element of
$K$. If either local component extended to a character on $G$, its
restriction to $\Gamma$ would be an abelian character nonzero on
$K=\overline{[\Gamma,\Gamma]}$, contrary to Kronecker--Weber.
Thus the common bit is meaningful with the multiplier-one
normalization but cannot be promoted in the suggested way to all
multipliers. The R2 element has lifts $(1,0),(0,1)$ and class one.
For an involution $(-1,f)$, R7 gives only
$\Delta((-1,f))=\delta((1,f))$; no value for an unclassified
involution follows. Conjugates of complex conjugation have value
zero.
\end{proof}

\begin{statement}{A.4}{Proved}{The phase generalization and its limit}
For a finite abelian phase group $M$, the quotient
$M^2/\operatorname{diag}M$ is canonically $M$ by
$(u,v)\mapsto uv^{-1}$. A pair of continuous characters on $\HH$
therefore defines the same gluing obstruction with values in $M$.
It descends to $G=\HH\Gamma$ if its ratio is conjugation invariant
and trivial on $K$. The supplied data do not specify a higher phase
character or a degree-two gerbe class.
\end{statement}
\begin{proof}
The kernel and surjectivity of the ratio map give the quotient.
The descent proof is the proof of A.2 with multiplication in $M$.
An equivariant $M$-torsor is governed by degree-one cohomology;
a gerbe would require a degree-two cocycle or a specified central
extension. Neither follows from a pair of signs. For roots of unity
with nontrivial cyclotomic action the coefficient module must also
be specified; ordinary characters then have to be replaced by the
appropriate crossed homomorphisms. We introduce no such extra data.
\end{proof}

\begin{statement}{A.5}{Proved}{What the proposed product formula actually supplies}
Set $X=\mathbb P^1_{\Q}\setminus\{0,1728,\infty\}$ and
$u=1-1728/j\in\mathcal O(X)^\times$. Its Kummer class has geometric
boundary residues $(1,1,0)$ in $\Ftwo$, ordered as
$(0,1728,\infty)$. For a rational point $t\in X(\Q)$ its pullback
is the square class of $1-1728/t$. This class is not by itself a
comparison of the two CM frames.
\end{statement}
\begin{proof}
The divisor on the coarse projective line is
$[1728]-[0]$, and $u(\infty)=1$. The residue of the Kummer double
cover $v^2=u$ at a puncture is its order modulo two, giving the
displayed vector and its zero sum. Functoriality of Kummer theory
gives the value at a rational section. For an arbitrary continuous
section of $\pi_1(X/\Q)\to G_{\Q}$, restriction still gives an
element of $H^1(G_{\Q},\mu_2)$, but identifying it with a rational
evaluation requires additional information about that section.
These residues are fixed properties of $u$, independent of $g$.
Calling $(-1)^{\Delta(g)}$ the third local sign gives a true
identity by A.2; identifying it with an independently defined
cusp datum is an additional comparison, not Weil reciprocity
applied to the displayed divisor. The section-reciprocity claims
ANABELIAN D.27--D.28 are not used as proofs of that comparison.
\end{proof}

\subsection{The exact scalar compression}

\begin{statement}{A.6}{Proved}{The CRT encoding is an encoding of parities}
There is a continuous homomorphism on $\HH$
\[
 u(h)\equiv(-1)^{a_3(h)}\pmod3,\qquad
 u(h)\equiv(-1)^{a_2(h)}\pmod4.
\]
It takes values in $\{1,5,7,11\}\subset(\Z/12)^\times$ and
has image modulo the global sign equal to $\delta$.
The complete table is
\[
\begin{array}{cc|ccc}
a_3&a_2&u&e&\delta\\\hline
0&0&1&0&0\\0&1&7&9&1\\1&0&5&4&1\\1&1&11&1&0
\end{array}
\]
where $e$ is the unique idempotent satisfying $2e=1-u$ in $\Z/12$.
The notation $(1-u)/2$ denotes this unique idempotent solution,
not division by an invertible element. For a general multiplier
only $[u(g\sigma^{-1})]$ is choice independent.
\end{statement}
\begin{proof}
CRT identifies the units with $\{\pm1\}\times\{\pm1\}$.
Multiplication of signs gives the homomorphism law. The idempotents
of $\Z/12\simeq\Z/3\times\Z/4$ are precisely the four pairs of
zeros and ones, and direct substitution gives the table. The two
solutions of $2e=1-u$ without the idempotence requirement need not
coincide. Changing $\sigma$ multiplies $u$ by $1$ or $11$, by
A.3. The ring $\Z/3\times\Z/4$ encodes the two signs as one
involution, but does not identify the two original frame matrices.
\end{proof}

\begin{statement}{A.7}{Refuted}{The uniform centre assertion at prime powers}
The assertion that
$\Ob(N)=Z(\SL_2(\Z/N))/\{\pm I\}$ is trivial at every prime
power is false. In particular $\Ob(8)\simeq C_2$.
For $N=2^a\prod_{i=1}^r p_i^{e_i}$ with $p_i$ odd,
the centre is the group of scalar square roots of one. Its rank
over $\Ftwo$ is $r+c(a)$, where
$c(0)=c(1)=0$, $c(2)=1$, and $c(a)=2$ for $a\ge3$.
For $N>2$, quotienting by the global sign reduces that rank by one.
\end{statement}
\begin{proof}
A matrix commuting with
$\left(\begin{smallmatrix}1&1\\0&1\end{smallmatrix}\right)$
has lower-left entry zero and equal diagonal entries; commuting
also with its transpose forces the upper-right entry zero. A
central matrix is therefore $aI$ with $a^2=1$. Conversely these
matrices are central. At an odd prime power,
$(a-1)(a+1)=0$ and the factors are coprime to each other at $p$,
so $a=\pm1$. At $2^a$, the roots are one for $a=1$, two for
$a=2$, and $\{1,-1,1+2^{a-1},-1+2^{a-1}\}$ for $a\ge3$,
as follows by comparing the two consecutive even valuations.
CRT finishes the proof. Modulo eight, $3I$ is central with
determinant one and is neither $I$ nor $-I$. The independent
enumeration in \texttt{check\_centres\_joint.py} checks the full
elementary-generator centralizers for $N=8,12$, and the roots for
$N=24,60$. The latter obstruction groups both have order four.
\end{proof}

\begin{statement}{A.8}{Refuted}{The proposed all-level formula is not a defined character}
The formula in Lead 8 cannot be used as a definition on all of
$G$ without new data: actual frames centralize individual torsion
elements, not the whole congruence group; moreover an arbitrary
multiplier is not a scalar element of $\SL_2(\Z/N)$.
\end{statement}
\begin{proof}
Even when an automorphism induces the identity on a finite
quotient, its frame satisfies $q(A)\in C_Q(q(r))$, which does
not imply $q(A)\in Z(Q)$. The noncentral Gaussian frame supplied
by A.30 below is an actual Galois example. We can strengthen it
to an example whose normalized congruence action is the identity.
The full Galois frame image in LB.4 supplies an element with
$T=-I$ and $W=S$. Every positive pure $r$-frame belongs to $F'$
(as proved in LB.5), and in $\SL_2(\mathbf F_3)^{\rm ab}=C_3$
the identity $T=q_3(A)^{-1}U_0^{8\rho}$ gives
$[T]=2\rho[U_0]$. Since $[-I]=0$ and $[U_0]$ generates $C_3$,
we obtain $\rho=0\pmod3$. Thus the R3 shear is the identity
modulo twelve, and both finite torsion images are fixed.
Nevertheless $W=S$ is noncentral projectively: it fails to commute
with the order-three generator. The finite matrix certificate
in \texttt{check\_pointed\_theta.py} checks that failure.
This refutes even the conditional centrality assertion when the
induced congruence action is literally the identity.
At $N=5$ the scalar
$2I$ has determinant four, not one. The diagonal matrix
$\operatorname{diag}(2,3)$ instead has determinant one but is
not central. Thus neither interpretation supplies the factor
$\langle\lambda\rangle$ in the asserted centre quotient. Also
R2 realizes one bit, and cannot by itself prove surjectivity to
groups of arbitrary rank. A.6 is the strongest all-multiplier
encoding established here from the supplied two parities.
\end{proof}

\subsection{A finite quotient which definitely cannot detect the sign}

\begin{statement}{A.9}{Refuted}{The order-6912 candidate does not determine $\Delta$}
Let $U=F'F^3$ and let
\[
 J=\ker(q_{12})\cap(F'F^{12})\triangleleft\B,
 \qquad Q_{\rm joint}=\B/J.
\]
Here the second factor is the actual Fermat quotient of the
level-two subgroup, retaining its extension by $S_3$.
The R2 element and the identity induce the same automorphism of
$Q_{\rm joint}$ but have different $\Delta$ values. Consequently
$\Delta$ does not factor through this finite action, even though
$|Q_{\rm joint}|=6912$.
\end{statement}
\begin{proof}
Under the signed CRT identification
$q_{12}(F)=V_4\times\SL_2(\mathbf F_3)$, the image of $U$ is
$V_4\times Q_8$, a group of order 32. Indeed
$\SL_2(\mathbf F_3)'=Q_8$, cubes kill its quotient of order three,
and cubes generate the $V_4$ factor. In the Fermat quotient
$F/F'F^{12}\simeq(\Z/12)^2$, the image of $U$ is
$3(\Z/12)^2\simeq(\Z/4)^2$. Its image in the joint quotient
is therefore a subgroup of a product of two finite 2-groups,
and is nilpotent. By the definition of $R_{\rm nil}(U)$,
\[
 R_{\rm nil}(U)\subseteq J.
\]
The R2 automorphism is the identity on $\B/R_{\rm nil}(U)$,
so it is the identity on $\B/J$. It has $\Delta=1$, whereas
the identity has $\Delta=0$. This proves nonfactorization with
the exact comparison map in the required direction.

The independent program realizes the Fermat quotient by the
Bass--Serre graph cocycle in $(\Z/12)[S_3]\rtimes S_3$, with
$P(s)=0$, $P(r)=e_1-e_r$. It takes its joint image with $q_{12}$,
enumerates all 6912 elements, and normally generates the image
of $U$ using $x^3,y^3,[x,y]$. The image has order 128; its
congruence and Fermat projections have orders 32 and 16.
These numbers verify the finite input; no finite shadow is
asserted to be a full GT element by the computation.
\end{proof}

\begin{statement}{A.10}{Proved}{Existence of a detecting characteristic level}
There is a finite characteristic quotient of $\B$ whose
induced outer action determines $\Delta$. Every such quotient
must have nonnilpotent image of $U=F'F^3$. Neither continuity
nor this existence theorem identifies a smallest such quotient.
\end{statement}
\begin{proof}
A finitely generated profinite group has finitely many open
subgroups of any bounded index: their coset actions are among
finitely many maps to finite symmetric groups. Intersecting all
open normal subgroups of index at most $n$ gives a cofinal
sequence of open characteristic subgroups $K_n$. The maps
$\Out(\B)\to\Out(\B/K_n)$ separate points. Indeed, if an
automorphism is inner at every level, the compact sets of
possible conjugators at successively finer levels have the
finite intersection property and give a conjugator in $\B$.
The same compactness argument shows these kernels are a
neighbourhood basis of the identity in $\Out(\B)$.
The specified embedding of $G$ into this outer automorphism
group is a homeomorphism onto its compact image. Since
$\ker\Delta$ is open, one of these kernels meets $G$ inside it.
If a finite quotient has nilpotent image of $U$, its kernel
contains $R_{\rm nil}(U)$ and the R2 element acts identically;
it therefore cannot determine $\Delta$ by A.9's argument.
\end{proof}

\begin{statement}{A.11}{Proved}{The compactness certificate and its correct quantifiers}
Let $P:G\to\widehat\Gamma_{0,5}^{\rm pure}$ be the continuously
evaluated pentagon word, and let $L_n$ be cofinal open normal
subgroups of its target. Then
\[
 \Delta|_{\GT}=0
 \quad\Longleftrightarrow\quad
 \text{for some }n,\quad
 \{g\in G:\Delta(g)=1,\ P(g)\in L_n\}=\varnothing.
\]
The corresponding finite test must retain actual $G$ residues
and the same evaluation maps; arbitrary separately constructed
finite shadows do not satisfy this criterion automatically.
\end{statement}
\begin{proof}
The sets on the right are closed and nested in the compact
space $G$. Their intersection is
$\Delta^{-1}(1)\cap P^{-1}(1)$, which is empty precisely
when the restriction vanishes. If all the closed sets were
nonempty, compactness would give a point in their intersection.
Conversely an empty set at one level implies the intersection
is empty. Refine a detecting level from A.10 together with
the finitely many word evaluations to obtain a finite test.
A family of candidate residues which has not been shown to
arise as projections of a common compact space of gentle
automorphisms can contain extra points. Surviving such points
do not prove the existence of a compatible profinite point.
\end{proof}

\subsection{The common CM field}

\begin{statement}{A.12}{Proved}{The maximal abelian closures have exactly one common extra bit}
Inside a fixed algebraic closure of $\Q$, put
$K_4=\Q(i)$ and $K_3=\Q(\sqrt{-3})$. Then
\[
 K_4^{\rm ab}\cap K_3^{\rm ab}
       =\Q^{\rm ab}(3^{1/4}),\qquad
 [\Q^{\rm ab}(3^{1/4}):\Q^{\rm ab}]=2.
\]
This proves the field-intersection assertion of Lead 7. It
does not prove its claimed surjectivity for the adelic frames.
\end{statement}
\begin{proof}
Both maximal abelian closures are Galois over $\Q$, since
their quadratic base fields are Galois and the maximal
abelian extension is characteristic. Their intersection $L$
contains $\Q^{\rm ab}$. In $G_L=\Gal(L/\Q)$ the two
distinct subgroups $A_i=\Gal(L/K_i)$ are abelian of index
two. Their intersection $Z$ is central: it commutes with
both subgroups, and those subgroups generate $G_L$.
Moreover $G_L/Z\simeq C_2^2$. Choose coset representatives
$a,b$ which generate this quotient. Every commutator is
a power of $[a,b]$, which is central, and
$[a,b]^2=[a^2,b]=1$. Thus the closed derived subgroup
has order at most two. Its fixed field is exactly
$\Q^{\rm ab}$, so $[L:\Q^{\rm ab}]\le2$.

Let $t=3^{1/4}$ and $M=\Q(i,t)$. The polynomial $X^4-3$
is Eisenstein, and its splitting group is the dihedral group
of order eight: take $r(t)=it$, $r(i)=i$, and
$c(t)=t$, $c(i)=-i$. The subgroup fixing $K_4$ is
$\langle r\rangle\simeq C_4$; the subgroup fixing
$it^2=\sqrt{-3}$ is $\langle r^2,rc\rangle\simeq C_2^2$.
Hence $M/K_4$ and $M/K_3$ are abelian. Composing with
$\Q^{\rm ab}$ keeps them abelian over both quadratic fields,
so $\Q^{\rm ab}(t)\subseteq L$. But $t\notin\Q^{\rm ab}$:
otherwise the nonabelian Galois extension $M/\Q$ would be
contained in an abelian extension. Since $t^2=\sqrt3$ is
already abelian, the degree is two and the upper bound is
attained. Finally $(3/4)^{1/4}=t/\sqrt2$, and $\sqrt2$
is abelian, giving the same bit as in R6.
\end{proof}

\begin{statement}{A.13}{Proved}{The actual finite CM fields in the stated models}
For $E_3:y^2=x^3+1$ and $E_4:y^2=x^3-x$, with
$r_0=(3/4)^{1/4}$, one has
\[
 \Q(E_3[4])=\Q(i,r_0)=\Q(x(E_4[3])).
\]
The common bit on $G_{\Q^{\rm ab}}$ is the action on $r_0$.
\end{statement}
\begin{proof}
For a nonzero two-torsion point of $E_3$ whose coordinate is
$e\in\{-1,-\zeta_3,-\zeta_3^2\}$, the four halves have
$x=e(1\pm\sqrt3)$. For $e=-1$ their ordinate squares
are $-9\pm6\sqrt3$; a square root of the plus value is
$\sqrt3(\sqrt3-1)r_0$. The product of the two ordinate
roots is $3\sqrt{-3}$, up to sign. These expressions lie
in $\Q(i,r_0)$ and generate it: $i$ is obtained by the
Weil pairing on four-torsion, $\sqrt3$ by the abscissae,
and $r_0$ from the displayed ordinate.
The three-division polynomial for $E_4$ is
$3X^4-6X^2-1$. Choose roots $a,b$ with
$a^2=1+2/\sqrt3$, $b^2=1-2/\sqrt3$, $ab=i/\sqrt3$.
Then $w=a+ib$ satisfies $w^4=4/3$, and
\[
 a=(w+2/w)/2,\qquad b=(w-2/w)/(2i).
\]
The splitting field is exactly $\Q(i,1/w)=\Q(i,r_0)$.
These algebraic identities are independently checked in
\texttt{check\_root\_algebra.py}. This is a direct verification
of the unrefereed field statement, separate from the supplied
frame-to-torsion comparison in R4.
\end{proof}

\subsection{The descent endomorphism, carefully typed}

\begin{statement}{A.14}{Proved}{The outer descent endomorphism exists}
Distinguish the abstract free profinite group $F_0=\langle X,Y\rangle$
from the level-two subgroup $F\subset\B$. Fix the identifications
of the prompt and the epimorphism
\[
 \varphi:F_0\twoheadrightarrow\B,
 \qquad X\mapsto s,\quad Y\mapsto sr.
\]
There is a continuous group endomorphism $D:G\to G$ with the
same multiplier, obtained by descending an outer action along
$\varphi$, then restricting to $F$ with its $S_3$ action
normalized. The outer compatibility group is exactly
$C_P=\operatorname{Fix}(D_P)$. Here $D_P=D$ uses the literal
current marking. It must not be identified with the original
manuscript's compatibility subgroup without changing markings.
\end{statement}
\begin{proof}
Put $Z=(XY)^{-1}$. The kernel of $\varphi$ is the closed
normal closure of $X^2,Z^3$. A gentle automorphism carries
each of $X,Y,Z$ to a conjugate of its common unit power.
It therefore preserves this kernel, in both directions, and
descends to a peripheral outer automorphism of $\B$.
Changing the representative by an inner automorphism of
$F_0$ changes the descendant by an inner automorphism of
$\B$, so descent is well defined and respects composition.

The subgroup $F=\ker(\B\to S_3)$ is characteristic:
the epimorphisms to $S_3$ form a single $\Aut(S_3)$ orbit.
An automorphism of $\B$ induces an inner automorphism on
$S_3$. Choose a lift $b\in\B$ of that conjugating
permutation and compose with $\operatorname{Ad}(b^{-1})$.
The normalized action induces the identity on $S_3$.
Its restriction to $F$ fixes the three boundary classes
up to the common unit power and commutes with the deck
$S_3$ in $\Out(F)$: the discrepancy between the images
of a deck lift is an element of $F$ and hence acts inner.
Different choices of $b$ differ by $F$, so define the
same outer restriction. Normalized representatives compose,
proving the homomorphism law for restriction.

This restriction is inverse to the given extension
$G\to\Out^{\rm per}(\B)$. To see injectivity explicitly,
an automorphism restricting inner to $F$ can be adjusted
to fix $F$ pointwise; then $b^{-1}\alpha(b)$ centralizes
$F$ for every $b\in\B$. The centralizer is trivial:
its $S_3$ image acts trivially on the faithful standard
action on $F^{\rm ab}$ and is therefore trivial, after
which the centre of the free profinite group is trivial.
Thus $\alpha$ is the identity. Continuity follows from
descent and restriction in the finite-quotient topology.
Composing restriction with descent gives $D$.

Compatibility of the two outer actions on $\B$ says
exactly that extension of $g$ equals descent of $g$;
applying the inverse restriction proves the fixed-point
criterion. The original manuscript uses a different free basis.
The required conversion, and a Galois counterexample to the
literal-current fixed-point assertion, are proved in A4.M1--A4.M5.
Thus the arithmetic functoriality sentence in checkpoint 01's
version of this statement was incorrect in the current marking
and is superseded here. The relevant manuscript passage is
\S47, cote 140--3, folios 415--420, scans 229--240.
No unproved equivalence with the harmonic equations is built
into this algebraic construction.
\end{proof}

\begin{statement}{A.15}{Proved}{The elementary factorization of the modular map}
With the Legendre coordinate $t$,
\[
 w=4t(1-t),\qquad j=64(4-w)^3/w^2
\]
has composite $256(1-t+t^2)^3/[t^2(1-t)^2]$.
This algebraic identity alone does not identify the
fixed-point equation of A.14 with relations (III$'$),
(NT1) and (NT2).
\end{statement}
\begin{proof}
Substitute $4-w=4(1-t+t^2)$ and
$w^2=16t^2(1-t)^2$. The factorization is over $\Q$.
An equality of coarse rational functions does not supply
the choices of tangential basepoints, inertia generators,
and conjugating paths needed to equate the respective
automorphism equations. Those are additional data in
the comparison requested by A.4(a).
\end{proof}

\begin{statement}{A.16}{Refuted}{The naive integral-specialization reading}
The stated Weierstrass models do not give two smooth
elliptic sections over $\Z$ whose torsion frames can simply
be read at both 2 and 3. The coarse $j$ coordinates do
coincide at those primes; this is weaker than the proposed
comparison of profinite frames.
\end{statement}
\begin{proof}
The discriminants are $-432$ for $E_3$ and $64$ for $E_4$.
Modulo three, $E_3$ has equation $y^2=(x+1)^3$ and is
singular at $(-1,0)$, whereas $E_4$ is smooth. Modulo
two both displayed cubics are singular. These assertions
follow by differentiating their homogeneous equations
and are checked in \texttt{check\_root\_algebra.py}.
After potential good reduction the common $j=0$ is
supersingular at these primes, but such a statement
requires a change of model and a base extension.
Moreover prime-to-residue-characteristic specialization
does not recover the wild-primary torsion signs used here.
The identity $1728=2^6 3^3$ proves the collision on the
coarse $j$ line, and supplies no map identifying the
two framed representations. Over $\Z[1/6]$ these two
coarse sections are disjoint.
\end{proof}
